\section{Reference implementation}
\label{ch:implementation}

The chatbot realizes the architecture through a shared conversational backend, document preparation tools, structured-data services and two client interfaces. This chapter explains the request and data paths needed to interpret the evaluation. It describes inspected source revisions; source-defined topology and behavior are not evidence of what any deployed instance ran.

\subsection{Repository and service organization}

The central \texttt{chat-bot} repository contains the Next.js web application, FastAPI service and GPU service. Its container configuration also declares an embedding worker that shares the API image, PostgreSQL storage and Valkey, an open-source Redis-compatible key-value store. Separating application processing from GPU inference allows their placement and resource allocation to differ; throughput and availability were not measured.

The \texttt{documents} repository holds original files and their Markdown representations. The \texttt{discord-bot} repository adapts Discord interactions to the shared API. The \texttt{mcp} repository contains a local Python MCP server and a pinned faculty-service submodule, and dedicated timetable and thesis-listing repositories supply further adapters. The project spans seventeen repositories; only inspected call sites, not repository names, are treated as evidence of a connection between them.

The account describes the chatbot at revision \texttt{859e7dc}, the documents repository at \texttt{16aa8c1}, the MCP repository at \texttt{8444c8a} and the Discord adapter at \texttt{3ca32b3}.

\subsection{Document preparation and ingestion}

Original institutional PDF and DOCX files are kept alongside their Markdown conversions. PDF conversion extracts embedded text, removes repeated edge text and page numbers, recognizes Macedonian legal-article headings and reflows wrapped lines. DOCX conversion extracts paragraphs. For scanned or unusable PDFs, an optional vision transcription route requests verbatim Macedonian text with lists, tables, article headings and explicit unreadable markers. The repository requires reviewing conversions against their originals; the thesis treats this as an editorial requirement and does not assume every file passed it. A missing exception or misextracted table can survive ingestion and later support a plausible but incomplete answer.

An upload client submits each processed Markdown file to the API with a stable name derived from its filename, a title and original-file metadata, and a subsequent request computes embeddings for eligible chunks. A separate, explicitly invoked synchronization can remove database documents that no longer exist locally. Revised content replaces the whole document and its chunks; unchanged chunks are not retained individually.

Runtime chunking belongs to the chatbot API. It strips HTML comments, splits on legal-article headings where available and otherwise on Markdown headings, and preserves section labels and chunk order. It uses a 1,300-character split target, a 1,650-character threshold below which a unit stays unsplit, and 150-character overlap, so short articles can remain intact. A tokenizer-aware helper in the document tooling uses different values and is not the active algorithm.

Citation resolution prefers an explicit authority URL, then a source URL, then a link to the original file in the documents repository. Optional Cloudflare R2 uploads mirror original files as an archive, but answers do not cite or retrieve from them. The schema can hold date, validity, verification and authority fields, but their population in the evaluated corpus is not established.

\subsection{Embedding storage and retrieval execution}

Embedding text combines a document title, section label and chunk body in a canonical Macedonian format. The default model is locally served BGE-M3, whose model family supports multilingual retrieval~\cite{chen2024m3}. The GPU service produces normalized embeddings, and PostgreSQL with pgvector stores the 1,024-dimensional vectors under a cosine-distance index based on hierarchical navigable small world (HNSW) graphs~\cite{malkov2020hnsw}. Database triggers mark vectors outdated when titles or chunks change, the worker recomputes them asynchronously, revision-guarded writes prevent stale embeddings from overwriting newer ones, and search uses only current embeddings.

The evaluated corpus audit recorded 74 FAQs and 1,403 chunks from nineteen documents. Retrieval searches the FAQ and document corpora separately before merging candidates. In the raw-query path used by the experiment, the dense limit is 31 candidates per corpus before merging and any permitted lexical FAQ expansion. The GPU cross-encoder reranks candidates using BGE-reranker-v2-m3~\cite{baaiBgeRerankerV2M3}. Final selection applies the ten-candidate budget, FAQ priority and neighboring-chunk expansion.

Context floors and the source-display threshold serve different purposes. The controlled comparison keeps the FAQ floor at 0.10, varies the document floor and whether one highest-scoring candidate is retained when all candidates fail it, and leaves source visibility at 0.30. The application separately falls back to dense ordering when the reranker fails. That operational recovery is distinct from the experimental empty-selection fallback; the study treats reranker failure as an execution failure rather than substituting another retrieval procedure.

\subsection{Web request and generation flow}

The Next.js server acts as a backend-for-frontend. It resolves a signed-in Google or Microsoft session to a backend user, loads the persisted conversation, assembles upstream messages and converts the stream. Its request to FastAPI \texttt{/chat/} carries server authentication, user identity, a response identifier and optional telemetry identifiers. This user identity is separate from the server-side credentials for the faculty's Central Authentication Service (CAS) used by some catalog adapters.

FastAPI first performs admission checks that resolve provider credentials and model access and apply sponsored quota rules. Evidence retrieval and useful-link selection then run as distinct steps, and the prompt combines their outputs with the question, date, available history and interface instructions. The provider layer supports hosted and local model families; the OpenAI route uses the Responses API and supplies combined tools to a LangChain agent. Support for multiple providers does not imply identical behavior across them.

The API emits a common stream of token, thinking, status, error, metadata, source and completion events, which the web adapter converts to the client's representation. Valkey supports resumable web streams, cancellation controllers are local to the web process, and conversations and messages are user-scoped PostgreSQL state.

One built-in tool recommends thesis committee candidates from database records and is combined with remotely discovered tools for the agent; its output is advisory. Tool availability establishes a possible request path, not that a suitable tool is chosen or its result correctly used, and the controlled study does not observe such agent runs.

\subsection{Structured records through MCP}

The chatbot's MCP client loads configured servers over streamable HTTP or server-sent events (SSE), can attach a service API key, filters tool names, and caches tool definitions for one hour; a server that fails to load does not disable the others. The local Python MCP service exposes read-only lookups over courses, staff, rooms, timetables and thesis listings, with Macedonian normalization, transliteration, fuzzy matching and ambiguity reporting. Timetable data comes through a Worker adapter over EduPage, the faculty's timetabling platform, and thesis listings through adapters that fetch faculty pages, authenticating through CAS where required. Because several layers cache records and may serve stale snapshots after a refresh failure, currentness claims require checking a deployed instance. A thesis link returned by a lookup is also not evidence that the thesis PDF was ingested into document RAG. A separate Spring-based faculty API exposes database-backed read-only MCP tools behind configurable service-level API-key filtering.

\subsection{Discord delivery and conversation continuity}

Discord requests start through slash commands such as \texttt{/ask}, \texttt{/prompt} and \texttt{/chat query}, or a message context action. Continuations arrive as replies to tracked answers or inside tracked private threads; there is no general mention trigger, so an arbitrary message is not an assistant request. The adapter resolves the Discord identity to a chatbot user, validates the model selection, and posts recent messages with \texttt{interface=discord} to the shared API, consuming SSE rather than calling a provider directly.

Rendering edits or splits messages at the 2,000-character platform limit. The adapter handles token, thinking, status, reset, error and completion events; thinking is display-only and excluded from saved history. Interface instructions favor simple Markdown, lists and plain URLs. The backend's structured-source event is not part of the adapter's event union and is ignored, although links embedded in answer text remain visible. This is a concrete source-presentation difference from the web client, which renders source cards.

Discord continuity is bounded and process-local: the adapter tracks at most 1,000 message or thread keys and sends nine prior turns plus the current turn, each limited to 2,000 characters. Split messages share one history entry, and conversations are not persisted through the web chat-state endpoints. Like/dislike buttons record feedback keyed to the response, client and user, linking feedback to an answer without restoring a lost conversation or validating its content.

\subsection{The implementation boundary used by the study}

The controlled study retains the actual retrieval function and prompt builder while intercepting generation. Captured prompts are submitted to controlled Luna through Chat Completions, whereas the production provider route uses Responses and an agent. The study fixes the prompt date and omits history and query rewriting. It excludes MCP discovery and execution, built-in tools, request-admission side effects, conversation saving, production analytics initialization and real client rendering. Schema validation remains observable, including overlong input rejection.

This design compares components without exercising the full live request lifecycle. Matching retrieval and prompt hashes support the provenance of those components, but controlled Luna defaults are not equated with every production provider setting, and measured generation latency omits channel delivery, admission and live tool execution. The implementation is therefore the source of the tested retrieval and prompt behavior, while the remaining paths supply architectural context and identified subjects for future evaluation.

\subsection{Changes after the evaluated revision}
\label{sec:post-evaluation-changes}

Between 8 September and 15 October 2026, after the evaluated revision \texttt{859e7dc}, several changes were merged into the chatbot. None of them is evaluated in this thesis. The reported results describe the evaluated revision with raw queries, and the changes are listed here only because they target stages that the evaluation identifies as weak.

\begin{itemize}
  \item \textbf{Query transformation.} History contextualization, rewriting and hypothetical-document generation previously ran only when the user had a credential for the transform model; otherwise retrieval silently used the raw query, so short follow-ups such as ``yes'' or ``and at the other faculty?'' retrieved nothing. Transforms now run on a sponsored hosted model when the user has no key for the transform model, transform output budgets were raised so that reasoning models do not return empty transforms, and the model that actually ran is logged. This addresses history-dependent follow-ups.
  \item \textbf{Romanized Macedonian.} Queries written in Latin script additionally receive an embedded Cyrillic transliteration as a query variant, and unchanged variants are no longer reported as partial fallbacks.
  \item \textbf{Short FAQ queries.} A bounded recovery targets short Macedonian FAQ queries about locations, hours and contacts while preserving the thresholds and source visibility.
  \item \textbf{Source authority and applicability.} FAQ evidence is constrained by institution scope and explicit exclusions: evidence for a different or unknown institution is filtered, and strict abstention applies only when a constraint can actually remove evidence. Document provenance is exposed in the model context, reranker relevance order is preserved, and the prompt clarifies source applicability.
  \item \textbf{Unsupported claims.} Recommendations about system functions must be grounded in a source or an authorized tool result.
  \item \textbf{Fallback diagnostics.} Empty or entirely invalid reranker results are recovered through the existing dense fallback, and diagnostic reasons are recorded for query-transform and reranker fallbacks.
  \item \textbf{Corpus synchronization.} A guarded synchronization loads the RAG corpus from the documents repository using a versioned bundle schema.
\end{itemize}

Together these changes address the failure stages discussed in this thesis: lost context in follow-ups, romanized Macedonian input, source authority and applicability, and unsupported claims. Whether they change answer quality is an empirical question, and evaluating them with the same protocol is left for future work.
